\subsection{Lists and Counting Principles}

Counting represents finite constructions by explicit lists and converts disjoint alternatives or successive choices into cardinality identities.

\begin{definition} \label{def:finite-sets-and-standard-index-sets} \hlt{Finite Sets and Standard Index Sets}\\
Set $J_0\defeq\emptyset$ and $J_n\defeq\{1,\ldots,n\}$ for $n\in\N$. A set $A$ is \hlt{finite} if it admits a bijection $e:J_n\to A$ with $n\in\N_0$.
\end{definition}

\begin{proposition} \label{prop:finite-cardinality-well-defined} \hlt{Finite Cardinality Is Well Defined}\\
Let $A$ be a set and let $m,n\in\N_0$. If $e:J_m\to A$ and $g:J_n\to A$ are bijections, then $m=n$.\\
Thus each finite set $A$ admits a unique $n\in\N_0$ such that a bijection $J_n\to A$ exists. Define the \hlt{finite cardinality} of $A$ by $\#A\defeq n$.
\begin{proof}
The composite $g^{-1}\circ e:J_m\to J_n$ is bijective.\\
It therefore suffices to prove that a bijection between two standard index sets forces their indices to agree.\\
For $r\in\N$, let $P_r$ assert that a bijection $J_{r-1}\to J_k$ with $k\in\N_0$ forces $k=r-1$. It suffices to prove $P_1$ and $(\forall r\in\N)\,(P_r\Rightarrow P_{r+1})$.\\
A bijection $J_0\to J_k$ forces $J_k=\emptyset$, hence $k=0$; thus $P_1$ holds. Assume $P_r$ and let $\varphi:J_r\to J_k$ be bijective. Since $J_r\neq\emptyset$, $k\in\N$.\\
Put $q\defeq\varphi(r)$. The restriction $\left.\varphi\right|_{J_{r-1}}$ is a bijection $J_{r-1}\to J_k\setminus\{q\}$.\\
Define $\rho:J_{k-1}\to J_k\setminus\{q\}$ by $\rho(i)\defeq i$ when $i<q$ and $\rho(i)\defeq i+1$ when $i\geq q$.\\
Its branch images $J_{q-1}$ and $J_k\setminus J_q$ are disjoint and cover $J_k\setminus\{q\}$. Both branches are injective, so $\rho$ is bijective.\\
The composite $\rho^{-1}\circ\left.\varphi\right|_{J_{r-1}}:J_{r-1}\to J_{k-1}$ is therefore bijective. Thus $k-1=r-1$ by $P_r$. Hence $k=r$, and $P_{r+1}$ holds.\\
Consequently, $(\forall r\in\N)\,P_r$ by Theorem \ref{thm:mathematical-induction}. Finally, $n=m$ by $P_{m+1}$ with the bijection $g^{-1}\circ e$.\\
\hlt{Structural Commentary}: Induction deletes one matched point and canonically relists the remaining index set.\\
\hlt{Key Insight}: Deletion reduces both finite indices while preserving bijectivity.
\end{proof}
\end{proposition}

\begin{definition} \label{def:lists-words-and-product-sets} \hlt{Lists, Words, and Product Sets}\\
Let $p\in\N$ and let $S_1,\ldots,S_p$ be sets.
\begin{enumerate}[label=(\roman*)]
\setlength{\itemsep}{0pt}
\item A \hlt{list} of length $p$ in $S_1\times\cdots\times S_p$ is a tuple $s=(s_1,\ldots,s_p)$ in that product.
\item If $S_1=\cdots=S_p=S$, a list in $S^p$ is a \hlt{word} of length $p$ over the alphabet $S$.
\item A counting model is \hlt{ordered} when its outcomes are lists and coordinate positions are retained.\\
Thus lists $s=(s_1,\ldots,s_p)$ and $t=(t_1,\ldots,t_p)$ are distinct iff $(\exists i\in J_p)\,s_i\neq t_i$.
\end{enumerate}
\end{definition}

\begin{proposition} \label{prop:addition-subtraction-and-division-principles} \hlt{Addition, Subtraction, and Division Principles}
\begin{enumerate}[label=(\roman*)]
\setlength{\itemsep}{0pt}
\item \hlt{Addition}: Let $r\in\N$ and let $A_1,\ldots,A_r$ be pairwise disjoint finite sets. Then $\bigcup\limits_{i=1}^{r}A_i$ is finite and $\#\bigl(\bigcup\limits_{i=1}^{r}A_i\bigr)=\sum\limits_{i=1}^{r}\#A_i$.
\item \hlt{Subtraction}: If $A$ is finite and $B\subseteq A$, then $B$ and $A\setminus B$ are finite and $\#(A\setminus B)=\#A-\#B$.
\item \hlt{Division}: Let $f:Y\to X$ be a surjection between finite sets. Its fibers are finite.\\
If $d\in\N$ satisfies $(\forall x\in X)\,\#f^{-1}(\{x\})=d$, then $\#Y=d\#X$ and $\#X=\frac{\#Y}{d}$.
\end{enumerate}
\begin{proof}
\begin{enumerate}[label=(\roman*)]
\setlength{\itemsep}{0pt}
\item For $k\in\N$, let $P_k$ denote claim (i) for $k$ sets. If $k=1$, then $\bigcup\limits_{i=1}^{1}A_i=A_1$ and $\sum\limits_{i=1}^{1}\#A_i=\#A_1$, so $P_1$ holds.\\
Assume $P_k$ and let $A_1,\ldots,A_{k+1}$ be pairwise disjoint finite sets. Put $U\defeq\bigcup\limits_{i=1}^{k}A_i$. By $P_k$, $U$ is finite and $\#U=\sum\limits_{i=1}^{k}\#A_i$.\\
Put $q\defeq\#U$ and $n\defeq\#A_{k+1}$, and choose bijections $e:J_q\to U$ and $g:J_n\to A_{k+1}$.\\
Define $h:J_{q+n}\to U\cup A_{k+1}$ as follows. Set $h(t)\defeq e(t)$ when $t\in J_q$ and $h(t)\defeq g(t-q)$ when $t\in J_{q+n}\setminus J_q$.\\
The two domain blocks partition $J_{q+n}$, and their images are the disjoint sets $U$ and $A_{k+1}$.\\
Both restrictions of $h$ are bijective, so $h$ is bijective.\\
Therefore $U\cup A_{k+1}$ is finite with cardinality $q+n=\sum\limits_{i=1}^{k+1}\#A_i$, so $P_{k+1}$ holds. Hence $(\forall k\in\N)\,P_k$ by Theorem \ref{thm:mathematical-induction}.
\item For $r\in\N$, let $P_r$ assert claim (ii) whenever $\#A=r-1$ and $B\subseteq A$. If $r=1$, then $A=B=\emptyset$, so $P_1$ holds.\\
Assume $P_r$, let $\#A=r$, and choose a bijection $e:J_r\to A$. Put $a\defeq e(r)$ and $A'\defeq A\setminus\{a\}$.\\
The restriction $\left.e\right|_{J_{r-1}}:J_{r-1}\to A'$ is bijective, so $\#A'=r-1$. If $a\in B$, put $B'\defeq B\setminus\{a\}$.\\
By $P_r$, the sets $B'$ and $A'\setminus B'$ are finite and $\#(A'\setminus B')=\#A'-\#B'$.\\
By part (i) for $B=B'\cup\{a\}$, $B$ is finite with $\#B=\#B'+1$.\\
Since $A\setminus B=A'\setminus B'$, the set $A\setminus B$ is finite and $\#(A\setminus B)=(r-1)-(\#B-1)=\#A-\#B$.\\
If $a\notin B$, then $B$ and $A'\setminus B$ are finite with $\#(A'\setminus B)=\#A'-\#B$ by $P_r$.\\
By part (i) for $A\setminus B=(A'\setminus B)\cup\{a\}$, $\#(A\setminus B)=((r-1)-\#B)+1=\#A-\#B$.\\
Thus $P_{r+1}$ holds. Consequently, $(\forall r\in\N)\,P_r$ by Theorem \ref{thm:mathematical-induction}.\\
An arbitrary finite $A$ and subset $B\subseteq A$ satisfy $P_{\#A+1}$, proving claim (ii).
\item Put $m\defeq\#X$. If $m=0$, then $X=\emptyset$ and hence $Y=\emptyset$ because $f:Y\to X$; both identities follow.\\
Suppose $m\in\N$, choose a bijection $e:J_m\to X$, and set $F_i\defeq f^{-1}(\{e(i)\})$ for $i\in J_m$.\\
Each $F_i$ is finite by part (ii). The sets $F_1,\ldots,F_m$ are pairwise disjoint with union $Y$.\\
The equality $(\forall i\in J_m)\,\#F_i=d$ holds by the uniform-fiber hypothesis.\\
The identity $\#Y=\sum\limits_{i=1}^{m}\#F_i$ holds by part (i). Hence $\#Y=\sum\limits_{i=1}^{m}d=md=d\#X$. Since $d\neq0$, $\#X=\frac{\#Y}{d}$.
\end{enumerate}
\hlt{Structural Commentary}: Induction constructs union and complement listings; fiber decomposition supplies division.\\
\hlt{Key Insight}: Disjoint blocks convert set decompositions into arithmetic.
\end{proof}
\end{proposition}

\begin{theorem} \label{thm:multiplication-principle-finite-products} \hlt{Multiplication Principle}\\
Let $p\in\N$ and let $n_1,\ldots,n_p\in\N_0$.
\begin{enumerate}[label=(\roman*)]
\setlength{\itemsep}{0pt}
\item \hlt{Successive choices}: Set $H_0\defeq\{\emptyset\}$.
\item At stage $i\in J_p$, let $(C_i(h))_{h\in H_{i-1}}$ be a family of finite sets satisfying $(\forall h\in H_{i-1})\,\#C_i(h)=n_i$.\\
Set $H_i\defeq\{(h,c):h\in H_{i-1},\ c\in C_i(h)\}$. Then $H_p$ is finite and $\#H_p=\prod\limits_{i=1}^{p}n_i$.
\item \hlt{Finite products}: Let $S_1,\ldots,S_p$ be finite sets satisfying $(\forall i\in J_p)\,\#S_i=n_i$.\\
Then $S_1\times\cdots\times S_p$ is finite and $\#(S_1\times\cdots\times S_p)=\prod\limits_{i=1}^{p}n_i$.
\end{enumerate}
\begin{proof}
\begin{enumerate}[label=(\roman*)]
\setlength{\itemsep}{0pt}
\item For $r\in\N$, let $P_r$ denote claim (i) with $p=r$. The map $(\emptyset,c)\mapsto c$ is a bijection $H_1\to C_1(\emptyset)$.\\
Thus $H_1$ is finite with $\#H_1=n_1$, so $P_1$ holds. Assume $P_r$ and consider an $(r+1)$-stage model.\\
By $P_r$ for its first $r$ stages, $H_r$ is finite with $\#H_r=\prod\limits_{i=1}^{r}n_i$.\\
If $H_r=\emptyset$, then $H_{r+1}=\emptyset$ and $\#H_{r+1}=0=n_{r+1}\#H_r=\prod\limits_{i=1}^{r+1}n_i$, so $P_{r+1}$ holds.\\
Suppose $H_r\neq\emptyset$. Put $m\defeq\#H_r$; then $m\in\N$. Choose a bijection $e:J_m\to H_r$.\\
Define the family $(A_j)_{j\in J_m}$ by $A_j\defeq\{e(j)\}\times C_{r+1}(e(j))$. The sets $A_1,\ldots,A_m$ are pairwise disjoint and have union $H_{r+1}$.\\
The map $c\mapsto(e(j),c)$ is a bijection $C_{r+1}(e(j))\to A_j$. Thus $A_1,\ldots,A_m$ are finite and $(\forall j\in J_m)\,\#A_j=n_{r+1}$.\\
Therefore $\#H_{r+1}=\sum\limits_{j=1}^{m}\#A_j=mn_{r+1}=\prod\limits_{i=1}^{r+1}n_i$ by Proposition \ref{prop:addition-subtraction-and-division-principles}(i). Thus $P_{r+1}$ holds.\\
Consequently, $(\forall r\in\N)\,P_r$ by Theorem \ref{thm:mathematical-induction}; in particular, $P_p$ holds.
\item In (i), take the constant families $C_i(h)\defeq S_i$. Then $H_p$ is finite with $\#H_p=\prod\limits_{i=1}^{p}n_i$.\\
Given $s=(s_1,\ldots,s_p)\in S_1\times\cdots\times S_p$, set $h_0\defeq\emptyset$. At stage $i\in J_p$, set $h_i\defeq(h_{i-1},s_i)$, and define $\Psi(s)\defeq h_p$.\\
The map $\Psi:S_1\times\cdots\times S_p\to H_p$ is bijective. Indeed, successive first and second coordinates recover the unique tuple $s$.\\
Hence $S_1\times\cdots\times S_p$ is finite. Its cardinality is $\#(S_1\times\cdots\times S_p)=\#H_p=\prod\limits_{i=1}^{p}n_i$.
\end{enumerate}
\hlt{Structural Commentary}: Induction counts recursive histories; a coordinate bijection transfers the count to products.\\
\hlt{Key Insight}: Uniform extension fibers multiply the number of histories.
\end{proof}
\end{theorem}

\begin{example} \label{ex:cards-binary-and-dna-counting} \hlt{Finite Words by the Multiplication Principle}\\
The following slot counts multiply by Theorem \ref{thm:multiplication-principle-finite-products}.
\begin{enumerate}[label=(\roman*)]
\setlength{\itemsep}{0pt}
\item A card in a standard deck is a pair consisting of one of $4$ suits and one of $13$ ranks. Thus the deck has $4\cdot13=52$ cards.
\item For $B\defeq\{0,1\}$, $\#(B^8)=2^8=256$, so the binary alphabet has $256$ words of length $8$.
\item For $D\defeq\{\mathrm{A},\mathrm{C},\mathrm{G},\mathrm{T}\}$, $(\forall N\in\N)\,\#(D^N)=4^N$. At $N=3$, the DNA alphabet therefore has $64$ codons.
\end{enumerate}
\hlt{Key Insight}: Fixed-size coordinate choices multiply, so constant alphabets produce exponential word counts.
\end{example}
